\section{\bt{System Overview}{系统概览}}
\label{sec:overview}
\bi{\system\ places a shared memory layer between data agents and a SQL database, PostgreSQL in our prototype, to reduce repeated exploration and its token and time cost. Successful sessions contribute executable metric and feature definitions, table descriptions, and verified join paths. Later agents retrieve this knowledge and obtain SQL for their periods or seed tables. Numbered revisions, learning records, and shared check results support validation and maintenance within the layer. Figure~\ref{fig:architecture} shows the publication, reuse, and maintenance paths.}
{\system\ 在数据智能体与 SQL 数据库之间设置共享记忆层，减少重复探索及其 token 和时间开销；原型采用 PostgreSQL。成功会话贡献可执行指标与特征定义、表描述以及经验证的连接路径。后续智能体检索这些知识，获得针对自身期间或种子表的 SQL。带编号的版本、学习记录与共享检查结果支撑层内的验证与维护。图~\ref{fig:architecture} 展示发布、复用与维护路径。}

\begin{figure*}[t]
\centering
% Vector drawing generated by tools/figures/architecture.py; rerun it after changing the figure.
\ifbilingual
\includegraphics[width=\textwidth]{figures/architecture-zh}
\else
\includegraphics[width=\textwidth]{figures/architecture}
\fi
\caption{\bt{\system\ shares database knowledge to reduce repeated exploration and its token and time cost. Successful sessions contribute metric and feature definitions (1--3). On a hit (4), agents reuse them for new periods and tasks, bypassing schema inspection, value probes, and join tests. Check results, maintenance, and repairs are shared as well (5--7). Snapshot-bound execution and per-seed temporal bounds connect this knowledge to each query or prediction instance (8--9). Thin arrows are calls, thick arrows write shared state, and dashed arrows read the database; purple, blue, and orange mark the learning, reuse, and execution paths.}{\system\ 通过共享数据库知识减少重复探索，节约 token 和时间。成功会话贡献指标与特征定义（1--3）；检索命中后（4），智能体将其用于新的期间和任务，省去重复的模式检查、取值探查与连接测试。检查结果、维护与修复同样共享（5--7）；绑定快照的执行机制和逐样本时间边界将这些知识关联到每个查询或预测样本（8--9）。细箭头为调用，粗箭头写入共享记忆，虚箭头读取数据库；紫、蓝、橙分别标出学习、复用与执行路径。}}
\label{fig:architecture}
\Description{A successful agent session contributes reusable knowledge through extraction and checked publication. The central shared knowledge store contains metric and feature definitions, condition dependencies, and versioned check results. Two user agents ask about September and October while holding the same revenue definition at v3; a third retrieves a spending feature at v1. A blue reuse lane carries the hit's SQL from the lookup step and fans out to these agents. A compact operation diagram contrasts schema inspection, value probes, and join tests with a definition hit that goes directly from lookup to SQL; it illustrates avoided exploration, not measured latency or LLM call counts. The check result for c2 and the three definition references to it use matching green highlights, showing that one result serves several definitions. Maintenance and metric repair write back for all consumers. Only T2 has changed, so c3 has no current result and its dependent metric is pending, while other matching results remain reusable. Metric execution checks and runs on one snapshot, and feature compilation excludes events at or after each seed time. The Database slab contains tables, table versions v7, v13 changing to v14, and v4, and the execution snapshot. Literal version values use the v prefix. Thick arrows write shared state and dashed arrows read database state.}
\end{figure*}


\bi{\emph{Learn and publish.} Agents explore through common table, join, and SQL tools. For metric questions, they submit answers with the query references and arithmetic that produced them; for prediction tasks, they submit feature SQL and validation scores. From successful sessions, \system\ extracts structured candidates, validates them through deterministic checks, and publishes accepted revisions. The publication gate controls updates to shared definitions. Table descriptions and verified join paths are also shared to reduce repeated exploration.}
{\emph{学习与发布。}智能体通过共用的表、连接和 SQL 工具开展探索。对于指标问题，提交答案时附上产生答案的查询编号和算式；对于预测任务，则提交特征 SQL 和验证分数。\system\ 从成功会话中提取结构化候选，执行确定性检查，并发布通过检查的版本。共享定义的更新由发布检查统一管理。表描述和经验证的连接路径也会共享，以减少重复探索。}

\bi{\emph{Retrieve and execute.} For analytics, \code{find\_metric} supplies a valid revision and its SQL. When the agent declares the revision to \code{run\_sql}, snapshot-bound execution checks its conditions on the query's snapshot. Prediction reuse has two complementary tiers: task programs provide executable starting points for retraining, and feature definitions support composition and refinement across compatible tasks. \code{find\_program} returns the previous task SQL and SQL compiled from its valid definitions; \code{find\_features} retrieves and compiles definitions for the entity, and \code{compose\_features} compiles a selected subset. The agent validates candidate SQL on the current task and audits temporal leakage before adoption.}
{\emph{检索与执行。}对于分析查询，\code{find\_metric} 返回有效版本及其 SQL；智能体在 \code{run\_sql} 中声明版本后，快照绑定执行机制在查询快照上检查条件。预测任务采用互补的两层复用：任务方案为重训提供可执行起点，特征定义支持兼容任务间的组合与细化。\code{find\_program} 返回此前任务 SQL 及有效定义的编译 SQL；\code{find\_features} 检索并编译实体上的定义，\code{compose\_features} 编译选定子集。智能体在当前任务上验证候选 SQL，并在采用前审计时间泄漏。}

\bi{\emph{Validate and maintain.} After an update, \system\ rechecks conditions whose dependencies changed, sharing matching check results across definitions and agents. It repairs affected metric definitions or invalidates them, and revokes feature definitions whose conditions fail. Cheaper equivalent metric revisions are checked before publication and shared by all consumers, amortizing optimization as well as maintenance.}
{\emph{验证与维护。}数据更新后，\system\ 重新检查依赖发生变化的条件，在定义与智能体之间共享匹配的检查结果。系统修复受影响的指标定义或将其置为失效，并撤销条件不再成立的特征定义。开销更低的等价指标版本在发布前经过检查，随后向所有使用者共享，从而共同分摊优化与维护开销。}
